\chapter{РЕАЛИЗАЦИЯ ПОЛЬЗОВАТЕЛЬСКИХ ПРИЛОЖЕНИЙ}

\section{Проектирование кроссплатформенного интерфейса}

Пользовательские приложения ориентированы на инженерные сценарии.
В них требуется быстро получать свойства рабочего тела
по ограниченному набору входных параметров
и проверять корректность режима по термодинамическим диаграммам.
Поэтому интерфейс должен поддерживать
одиночные расчёты, пакетную обработку и визуализацию.

Базовые функциональные требования к интерфейсу:
\begin{itemize}
\item выбор режима ввода (например, $pT$, $ph$, $ps$, $px$, $\rho T$) и отображение полного набора свойств;  
\item явное указание единиц измерения и возможность контроля диапазонов;  
\item диагностируемая обработка ошибок и отображение причин отказа в расчёте;  
\item табличный режим с построчной локализацией ошибок;  
\item построение диаграмм (например, $T$\textendash$s$, $p$\textendash$h$, $p$\textendash$v$) и отображение купола насыщения;  
\item экспорт результатов (таблицы, графики, журналы) в распространённые форматы.  
\end{itemize}

Нефункциональные требования:
\begin{itemize}
\item кроссплатформенность и воспроизводимая сборка;  
\item отзывчивость интерфейса при интерактивной работе и при пакетных расчётах;  
\item отсутствие дублирования предметной логики в интерфейсном слое;  
\item единые контракты обмена между клиентской и серверной частями через \texttt{if97\_app\_api}.  
\end{itemize}

	Для выполнения этих требований в проекте реализованы две интерфейсные оболочки,
	основанные на едином вычислительном ядре \texttt{if97\_core} и общих DTO \texttt{if97\_app\_api}.  
	Первая оболочка представляет собой нативное настольное приложение на FLTK~\cite{fltk_rs_docs}.  
	Вторая оболочка построена на связке Yew, WebAssembly и Tauri,
	что обеспечивает кроссплатформенный пользовательский интерфейс и доступ к функциям ОС
	~\cite{yew_docs,webassembly_docs,tauri_docs}.  

	Ключевые экраны настольной части на базе Tauri приведены на рисунках ниже:
	одиночный расчёт, табличная обработка, построение диаграмм и служебный журнал.

\section{Наследуемый интерфейс: настольное приложение на FLTK}

	Нативный интерфейс реализован на базе библиотеки FLTK через Rust-биндинги \texttt{fltk-rs}~\cite{fltk_rs_docs}.  
Данный вариант ориентирован на локальное использование без веб-окружения исполнения и сетевой инфраструктуры,
что важно для инженерного рабочего места и демонстрации автономности решения.  

Пользовательская часть организована в виде набора вкладок,
соответствующих типовым сценариям:
«Одиночный расчёт», «Табличный расчёт», «Графики», «О программе».  
Одиночный режим обеспечивает быстрый расчёт одной точки по выбранной входной паре.  
Табличный режим предназначен для пакетной обработки нескольких строк входных данных.  
Графический режим предназначен для визуального анализа режима на диаграммах и для подготовки иллюстративных материалов.  

	С архитектурной точки зрения FLTK-клиент реализован как тонкий адаптерный слой над ядром.  
Вся предметная логика остаётся в \texttt{if97\_core}, а интерфейс выполняет ввод, валидацию формата, вызов ядра и представление результата.  
	Для повышения надёжности реализации применён подход, близкий к \texttt{MVU} (\texttt{Model--View--Update}),
	где обработка событий и изменения состояния происходят через обмен сообщениями.  
Взаимодействие между потоками пользовательского интерфейса и логикой приложения строится через \texttt{fltk::app::channel},
что позволило уйти от схемы с \texttt{Arc<Mutex<T>>},
уменьшить связность модулей и снизить риск взаимных блокировок.  

	Важная особенность FLTK-клиента --- поддержка наборов данных.  
Это позволяет сохранять и переиспользовать рассчитанные точки и таблицы,
а также управлять их видимостью на диаграммах.  
Такой подход удобен для сравнения нескольких серий расчётов, разных режимов или экспериментальных данных.  

	Визуализация в FLTK выполняется через \texttt{plotters}.
Рендер строится в буфер, после чего выводится в интерфейс.
Это уменьшает мерцание при обновлении графиков.
Для устойчивости отображения применяется фильтрация \texttt{NaN}
и отсечение линий до передачи данных в рендерер.

\section{Современный стек: интеграция Yew, WASM и Tauri}

	Вторая оболочка реализована на связке веб-технологий и настольной упаковки.
	Клиентская часть построена на фреймворке Yew и компилируется в WebAssembly
	(сборка через Trunk).
	Настольные функции, вычисления и доступ к файловой системе
	реализованы в серверной части Tauri
	~\cite{yew_docs,webassembly_docs,tauri_docs}.  

Принципиальный момент архитектуры состоит в том,
что клиентская часть не выполняет термодинамические расчёты самостоятельно.  
	Вместо этого интерфейс формирует запросы в формате DTO (слой \texttt{if97\_app\_api}) и вызывает команды Tauri через механизм \texttt{invoke}.  
	Серверная часть Tauri, в свою очередь, выполняет вычисления через \texttt{if97\_core} и возвращает сериализуемый результат обратно в интерфейс.  
Это обеспечивает единый источник истины по вычислениям и совпадение результатов с другими каналами доставки.  

На рисунке~\ref{fig:tauri-single} показан экран одиночного расчёта.
На этом экране пользователь выбирает входную пару, задаёт параметры состояния и сразу получает полный набор термодинамических свойств.

\begin{figure}[htbp]
    \centering
    \includegraphics[width=\textwidth]{tauri_single.png}
    \caption{Окно одиночного расчёта в приложении на базе Tauri}
    \label{fig:tauri-single}
\end{figure}

	Для обмена данными важна совместимость с сериализацией в \texttt{JSON} и типичными ограничениями веб-окружения.  
	Поэтому DTO-слой учитывает необходимость корректного кодирования специальных числовых значений (например, \texttt{NaN} и \texttt{Infinity})
	и хранения диагностических сообщений об ошибках.  
	Это дополнительно проверяется тестами \texttt{if97\_app\_api},
	которые подтверждают сериализацию и обратное чтение
	нечисловых значений как строковых маркеров,
	пригодных для \texttt{JSON}-обмена с клиентской частью.

Термодинамические расчёты относятся к вычислительно интенсивным задачам.  
Если запускать их внутри асинхронной среды исполнения без изоляции,
возможны зависания и деградация отзывчивости интерфейса.  
	Поэтому тяжёлые вычисления в серверной части Tauri выполняются в блокирующем пуле через \texttt{spawn\_blocking},
что отделяет вычислительную нагрузку от интерфейса и операций ввода-вывода
~\cite{tokio_spawn_blocking_docs}.  

На уровне архитектуры интерфейса решается задача сохранения введённых данных и результатов при переключении вкладок.  
Для этого применяется принцип «единого источника истины»: состояние вынесено в общий контекст приложения,
в том числе через \texttt{AppContext} и \texttt{PersistentState}.
	Навигация реализована через управление видимостью (классы \texttt{CSS}),
	без разрушения элементов \texttt{DOM} и потери пользовательского ввода.  

\section{Визуализация: построение диаграмм и экспорт данных}

Обе интерфейсные реализации поддерживают два базовых режима работы: одиночный и табличный.  
Одиночный режим используется для интерактивного ввода одной точки и анализа полного набора свойств состояния.  
Табличный режим ориентирован на пакетные инженерные сценарии,
когда требуется обработать массив входных данных и получить итоговую таблицу результатов.  

Ключевое требование табличного режима --- устойчивость к некорректным строкам.  
В реальных данных отдельные строки могут выходить за пределы области применимости IF97 или содержать ошибки ввода.  
	Поэтому таблица обрабатывается построчно,
	а ошибки локализуются на уровне конкретной строки без остановки всего расчёта.  

На рисунке~\ref{fig:tauri-table} показан табличный режим.
Он удобен для массовой проверки исходных данных, где отдельные строки могут выходить за пределы области применимости или содержать ошибки ввода.

\begin{figure}[htbp]
    \centering
    \includegraphics[width=\textwidth]{tauri_table.png}
    \caption{Окно табличного расчёта с построчной обработкой ошибок}
    \label{fig:tauri-table}
\end{figure}

Визуализация результатов выполняется через построение термодинамических диаграмм.  
Диаграммы используются для контроля режима и для наглядного представления результатов в отчёте.  
Для корректной интерпретации состояния на диаграмме важна отрисовка купола насыщения,
то есть границ $x=0$ (насыщенная жидкость) и $x=1$ (насыщенный пар).  

На рисунке~\ref{fig:tauri-plot} приведён экран построения диаграммы.
По нему видно, как интерфейс отображает экспериментальные точки, рассчитанные кривые и купол насыщения в одном рабочем окне.

\begin{figure}[htbp]
    \centering
    \includegraphics[width=\textwidth]{tauri_plot.png}
    \caption{Построение термодинамической диаграммы с куполом насыщения}
    \label{fig:tauri-plot}
\end{figure}

Построение купола насыщения выполняется с учётом требований к гладкости и стоимости вычислений.  
Для диаграмм, где по оси откладывается удельный объём $v$,
применяются меры для повышения плотности точек в области малых значений $v$,
в частности использование метрики, основанной на $\ln(v)$.  
	В оболочке на базе Yew, WebAssembly и Tauri дополнительно используется
	контроль актуальности асинхронных ответов по \texttt{req\_id}.
Купол насыщения кэшируется по ключу \texttt{(chart\_type, swap\_axes)}.
Это устраняет «вспышки» устаревших кривых
при частом изменении параметров отображения.  

Экспорт данных необходим как для инженерной работы, так и для подготовки материалов к защите.  
В проекте поддерживаются следующие типовые варианты экспорта:
\begin{itemize}
\item сохранение табличных результатов в CSV;  
\item экспорт диаграмм в растровые изображения (PNG) для вставки в текст ВКР;  
\item выгрузка журналов и диагностических сообщений для анализа ошибок и воспроизводимости результатов.  
\end{itemize}

	В FLTK-версии экспорт реализуется как сохранение подготовленных файлов из нативного приложения.  
	В Tauri-версии экспорт выполняется через команды серверной части и системные диалоги,
	что обеспечивает корректную интеграцию с ОС и единый пользовательский опыт на разных платформах~\cite{tauri_docs}.  
	Для диагностики и контроля состояния приложения предусмотрен раздел «О приложении» и системный журнал,
	позволяющий проверять версию сборки и анализировать выполнение вычислений.

На рисунке~\ref{fig:tauri-log} показаны эти служебные элементы интерфейса.

\begin{figure}[htbp]
    \centering
    \includegraphics[width=\textwidth]{tauri_log.png}
    \caption{Раздел «О приложении» и системный журнал}
    \label{fig:tauri-log}
\end{figure}

	Для серверной части Tauri
	в \path{wasm/src-tauri/src/lib.rs}
поддерживаются проверочные тесты
чистых функций расчёта, парсинга табличного ввода
и генерации купола насыщения.
Они не зависят от графических диалогов,
поэтому пригодны для автоматического прогона в системе непрерывной интеграции.

\section[Сравнительный анализ интерфейсных решений]{Сравнительный анализ производительности и сценариев применения интерфейсных решений}

Наличие двух независимых оболочек демонстрирует переносимость вычислительного ядра и корректность архитектурного решения «одно ядро --- несколько потребителей».  
Обе реализации используют \texttt{if97\_core}, поэтому при одинаковом вводе обеспечивают совпадение численных результатов.  
Различия относятся к технологическому стеку, способам доставки и эксплуатационным сценариям.  

	FLTK-версия является лёгкой и автономной.  
Она подходит для быстрого запуска, офлайн-использования и работы на маломощных ПК,
а также удобна для демонстрации минимального слоя между пользователем и ядром.  

	Версия на базе Yew, WebAssembly и Tauri ориентирована на современный пользовательский интерфейс,
	кроссплатформенную упаковку и интеграцию с системными возможностями (файловые диалоги, экспорт, журналы)~\cite{tauri_docs}.  
Эта реализация удобна для расширения интерфейса, визуализации и сценариев,
где важны единые паттерны пользовательского взаимодействия и переносимость веб-подходов на настольную среду.  

С точки зрения эксплуатационного профиля
	FLTK-решение выигрывает простотой запуска,
	меньшим числом промежуточных слоёв
	и предсказуемым поведением на локальной рабочей станции.
	Связка Yew, WebAssembly и Tauri, напротив,
	даёт более богатый интерфейсный слой
	и лучшую основу для дальнейшего развития пользовательского опыта,
	но требует более сложного стека сборки
	и аккуратного разделения клиентской и серверной частей.

Таким образом, две оболочки решают одинаковую предметную задачу,
но оптимальны для разных условий применения,
а единый вычислительный слой обеспечивает корректность и воспроизводимость результатов.  
